\section{Cycle--Transport Rigidity}

Let (G=(V,E)) be a finite graph with maximum degree
(\maxdeg(G)\le \Delta).
Fix (R\in\mathbb N).

Let
[
\mathcal C_R(G)
]
denote the set of cycles contained in radius--(R) balls and define the
cycle--overlap rank
[
COR_R(G)=\dim_{\mathbb F_2}\mathrm{span}{\chi_C:C\in\mathcal C_R(G)}
\subseteq \mathbb F_2^{E}.
]

\subsection{Edge--Cycle Incidence}

Let (C_1,\dots,C_m) be a basis of
(\mathrm{Im},\Phi_R(G)) where
[
m = COR_R(G).
]

Define the edge--cycle incidence matrix
[
A\in\mathbb F_2^{E\times m}
]
by
[
A_{e,j}=
\begin{cases}
1 & e\in C_j,\
0 & e\notin C_j.
\end{cases}
]

Then
[
\operatorname{rank}(A)=m.
]

\subsection{Cycle--Transport Signature}

For a vertex (v\in V) define the transport vector

[
\tau_R(v,u)
===========

#{C_j:(v,u)\in C_j}
\qquad u\in N(v).
]

The \emph{cycle--transport signature} is

[
\Theta_R(v)
===========

\Big(
\deg(v),
{\tau_R(v,u)}_{u\in N(v)}
\Big).
]

\subsection{Homogeneity Assumption}

Assume
[
\Theta_R(v)=\Theta_R(w)
\qquad
\forall v,w\in V.
]

Under bounded degree,
the number of possible radius--(R) neighbourhood types is finite.
Hence there exists a constant

[
T_0 = T_0(\Delta,R)
]

such that the rows of (A) fall into at most (T_0) equivalence classes.

\subsection{Rank Constraint}

A matrix with at most (T_0) distinct row types satisfies

[
\operatorname{rank}(A)\le T_0.
]

But

[
\operatorname{rank}(A)=COR_R(G)=m.
]

Thus

[
m \le T_0(\Delta,R).
]

\subsection{Rigidity Lemma}

\begin{lemma}[Cycle--Transport Rigidity]
There exists a constant (T=T(\Delta,R)) such that if
[
COR_R(G)\ge T
]
then
[
\exists u,v\in V(G):
\Theta_R(u)\ne\Theta_R(v).
]
\end{lemma}

\begin{proof}
If all transport signatures were identical,
the edge--cycle incidence matrix would have at most
(T_0(\Delta,R)) row types, forcing
(\operatorname{rank}(A)\le T_0(\Delta,R)).
For
(COR_R(G) > T_0(\Delta,R))
this contradicts
(\operatorname{rank}(A)=COR_R(G)).
\end{proof}

\subsection{WL(^2) Consequence}

The (2)-dimensional Weisfeiler--Leman refinement distinguishes vertices
with different neighbour transport profiles. Hence

[
\Theta_R(u)\ne\Theta_R(v)
\Rightarrow
\chi_{WL^2}(u)\ne\chi_{WL^2}(v).
]

\subsection{FO(^3) Consequence}

WL(^2) equivalence coincides with the (3)-pebble
Ehrenfeucht--Fra"{\i}ss'e equivalence on graphs.
Therefore

[
COR_R(G)\ge T
\Rightarrow
\exists u,v:\ tp_3(u)\ne tp_3(v).
]

\subsection{Structural Chain}

[
COR_R(G)
\Rightarrow
\Theta_R\text{ diversity}
\Rightarrow
WL^2
\Rightarrow
FO^3.
]

This establishes cycle--overlap rank as a structural source of
first--order distinguishability.

